% Estilo de listagem para C++ (competitivo) — rosa, cinza e roxo claro
\definecolor{codebg}{RGB}{248,244,252}        % fundo roxo bem claro
\definecolor{codeframe}{RGB}{195,188,210}      % borda cinza-arroxeado
\definecolor{codetext}{RGB}{62,58,72}         % texto base cinza escuro
\definecolor{codekeyword}{RGB}{210,88,138}     % palavras-chave rosa
\definecolor{codestring}{RGB}{72,130,85}       % strings verde
\definecolor{codecomment}{RGB}{165,140,195}   % comentários roxo claro
\definecolor{codenum}{RGB}{165,158,180}        % numeração de linhas

\lstdefinelanguage{cppmaratona}{
  language=C++,
  morekeywords={ll,ull,ld,vi,vll,pii,pll,iii,lll,vi,vc,pb,mp,fi,se,eb,mt,all},
  sensitive=true
}

\lstdefinestyle{maratona-cpp-base}{
  language=cppmaratona,
  extendedchars=true,
  backgroundcolor=\color{codebg},
  basicstyle=\ttfamily\small\color{codetext},
  keywordstyle=\color{codekeyword}\bfseries,
  stringstyle=\color{codestring},
  commentstyle=\color{codecomment}\itshape,
  numberstyle=\tiny\color{codenum},
  numbers=left,
  numbersep=8pt,
  stepnumber=1,
  tabsize=2,
  showspaces=false,
  showstringspaces=false,
  breaklines=true,
  breakatwhitespace=false,
  frame=single,
  rulecolor=\color{codeframe},
  xleftmargin=12pt,
  xrightmargin=6pt,
  framexleftmargin=12pt,
  captionpos=b,
  aboveskip=1em,
  belowskip=1em
}

% Acentos em comentários (precisa de \begin{document} para \c)
\AtBeginDocument{%
  \lstdefinestyle{maratona-cpp}{%
    style=maratona-cpp-base,
    literate=
      {á}{{\'a}}1 {â}{{\^a}}1 {ã}{{\~a}}1 {à}{{\`a}}1
      {é}{{\'e}}1 {ê}{{\^e}}1 {í}{{\'i}}1
      {ó}{{\'o}}1 {ô}{{\^o}}1 {õ}{{\~o}}1 {ú}{{\'u}}1
      {ç}{{\c c}}1 {Ç}{{\c C}}1
  }%
}

\lstset{style=maratona-cpp-base}

% Diretório raiz dos códigos (relativo à pasta apostila/)
\newcommand{\codigosdir}{../codigos}

% Inclui um arquivo .cpp: \incluircpp{mineira_2024/copos_de_leite.cpp}
% Opções extras do listings: \incluircpp[firstline=5,lastline=20]{...}
\newcommand{\incluircpp}[2][]{%
  \lstinputlisting[style=maratona-cpp,#1]{\codigosdir/#2}%
}

% Com legenda e referência cruzada (\ref{lst:...})
\newcommand{\incluircppcaption}[3][]{%
  \lstinputlisting[
    style=maratona-cpp,
    caption={#3},
    label={lst:#2},
    #1
  ]{\codigosdir/#2}%
}
